% NSF MSPRF Project Summary (NSF 23-603)
%
% Limit: 1 page. Required: separate headings "Overview", "Intellectual Merit", and
% "Broader Impacts" (a missing heading means return without review). Third person,
% understandable to a scientifically literate lay reader; not an abstract of the proposal.
% Uploading as a PDF is allowed only because the summary needs special characters.
%
% PAPPG formatting: Computer Modern at >= 11pt, margins >= 1in. Uploaded sections should
% leave out page numbers (Research.gov paginates). Nothing may be promised here that is
% not in the Project Description.

\documentclass[11pt]{article}

\usepackage[margin=1in]{geometry}
\usepackage{amsmath, amssymb}
\usepackage{microtype}
\usepackage{enumitem}
\usepackage[compact]{titlesec}

\titleformat{\section}{\large\bfseries}{\thesection.}{0.5em}{}
\titlespacing*{\section}{0pt}{1.2ex plus .4ex minus .2ex}{0.5ex plus .2ex}
\setlist{nosep, leftmargin=*}

% Block paragraphs, as in the Project Description.
\setlength{\parindent}{0pt}
\setlength{\parskip}{0.3\baselineskip plus 2pt minus 1pt}

\newcommand{\Gm}{\mathbb{G}_m}
\newcommand{\Sp}{\mathrm{Sp}}
\newcommand{\SL}{\mathrm{SL}}

\pagestyle{empty}

\begin{document}

\section*{Overview}

The PI will study the Hodge theory of Lagrangian fibrations at MIT, with sponsoring
scientist Davesh Maulik. Hyper-K\"ahler manifolds are basic building blocks of compact
K\"ahler manifolds with trivial canonical bundle, and classifying them is a fundamental open
problem. By the recently proven hyper-K\"ahler SYZ conjecture, every known deformation type
of hyper-K\"ahler manifolds admits Lagrangian fibrations, whose general fibers are 
abelian varieties of half the dimension. The project asks which abelian varieties 
occur as such fibers, i.e.\ where the associated period map lands.

The PI's thesis proves that for a non-isotrivial Lagrangian fibration of a $2n$-dimensional
hyper-K\"ahler manifold, the monodromy is irreducible and the very general fiber has
Mumford--Tate group $\Gm\cdot(\Sp_{2h})^{n/h}$ up to isogeny, for some $h\mid n$. 
Equivalently, the fiber is isogenous to a product of $h$-dimensional abelian varieties, 
each \emph{Hodge-generic} (with no unexpected Hodge classes). Moreover, $h=n$ for a 
general fibered member of each deformation type with $b_2\ge7$. 
The research objectives are:
\begin{enumerate}[label=\textbf{RO \arabic*.}]
  \item Determine which $h$ occur, and whether $h<n$ forces $K3^{[n]}$ type, using the
  Donagi--Markman cubic condition to split the period map along algebraic foliations of the
  base.
  \item Study non-compact Lagrangian fibrations, chiefly the $\SL_r$ Hitchin fibration (whose
  monodromy the PI's methods show to be irreducible), and determine the Mumford--Tate group
  of its very general fiber via Lefschetz theory for spectral curves.
  \item Work toward the hyper-K\"ahler Bogomolov--Mumford conjecture (a uniruled divisor in
  every ample linear system), using typical intersections in the moduli of abelian varieties
  to find families of fibers with an elliptic factor, whose degenerations give rational
  curves.
\end{enumerate}

Maulik, a leading expert on abelian fibrations, hyper-K\"ahler varieties, and Hitchin
systems, will meet with the PI weekly; at MIT the PI will learn the sheaf-theoretic tools
the project needs, such as perverse filtrations, Fourier--Mukai transforms, and moduli of
sheaves on K3 surfaces.

\section*{Intellectual Merit}

Since the SYZ theorem makes Lagrangian fibrations the gateway to classifying hyper-K\"ahler
manifolds, constraints on their fibers bear directly on classification. RO~1 would show that
decomposable Lagrangian fibrations are exactly those built from elliptic fibrations, as for
Hilbert schemes of points on elliptic K3 surfaces, and would characterize the $K3^{[n]}$
deformation type Hodge-theoretically. RO~2 would give a new Hodge-theoretic invariant of the
Hitchin system, a fundamental integrable system in non-abelian Hodge theory and geometric
Langlands, complementing recent work on its topology and algebraic cycles. RO~3 would bring
functional transcendence to bear on rational curves on hyper-K\"ahler manifolds, which are
tied to their algebraic cycles.

\section*{Broader Impacts}

The PI will organize a learning seminar on Lagrangian fibrations and Hitchin systems for
Boston-area graduate students and postdocs, and seek support for an early-career workshop.
Both will be open to all and aimed at researchers without local access to the area, with
notes and recorded talks posted publicly. The PI will also volunteer with Boston-area math
circles and competitions such as HMMT, favoring programs serving schools with limited
access to advanced mathematics, and seek to mentor undergraduates through the MIT Summer
Research Program. MIT's teacher-training program for new postdocs will further develop the
PI as an educator. Finally, Lagrangian fibrations and Hitchin systems are integrable systems
that appear throughout mirror symmetry and mathematical physics, so the results will be
useful beyond algebraic geometry.

\end{document}
